\documentclass[10pt,a4paper]{amsart}
\usepackage[margin=19mm]{geometry}
\usepackage{amsmath,amssymb,mathtools}
\usepackage[scr=rsfs,cal=euler]{mathalfa}
\usepackage{xfrac}
\usepackage[unicode,hidelinks,bookmarks=false]{hyperref}
\newcommand{\ie}{i.e.\ }
\newcommand{\cf}{cf.\ }
\newcommand{\resp}{respectively\ }
\newcommand{\Subsection}{\subsection}
\providecommand{\nobreakdash}{\nobreak\hbox{-}}
\setlength{\emergencystretch}{2em}
\multlinegap0pt
\usepackage{amssymb}
\newtheorem{theorem}{Theorem}
\newtheorem{corollary}{Corollary}
\title{The transfer of Artinian and Cohen-Macaulay properties under module-finite extensions}
\author{Tran Do Minh Chau}
\date{Source: arXiv:2609.11003v1 (CC BY 4.0)}
\begin{document}
\maketitle

\noindent\emph{Source and licence.} The transfer of Artinian and Cohen-Macaulay properties under module-finite extensions. Version arXiv:2609.11003v1 (CC BY 4.0), distributed by arXiv under the Creative Commons Attribution 4.0 International licence, http://creativecommons.org/licenses/by/4.0/, as recorded in the arXiv licence metadata for this exact version and shown on the abstract page https://arxiv.org/abs/2609.11003v1. Full licence notice: \texttt{licenses/arxiv-2609.11003-CC-BY-4.0.txt}.\par\smallskip

\noindent\textit{Editorial scope.} Counted unit: the article's corollary C:2 (source0.tex lines 331--332). The article's complete original proof of it is delivered with it (source0.tex lines 334--341, one \texttt{proof} environment of 8 lines). No mathematics is written, repaired or re-derived: the statement and the proof are the article's own lines, in the article's own notation, and no retained line is substituted or re-flowed.  Retained uncounted as the source-local context the proof needs: lines 284--291.  Nothing was omitted from any retained span, so the package carries no omission record.  No retained line contains a citation, so the wrapper carries no bibliography and no bibliography entry.  No retained line contains a figure environment, an image-inclusion command or drawing code, so no image is delivered and the package declares no asset dependency.  Editorial formatting only, in addition: an AMS \texttt{amsart} preamble that restates the article's own declarations and macros used by the retained lines, namely the environments \texttt{theorem}, \texttt{corollary} on separate counters, together with the standard packages those lines need.  The retained environments print in this document as Theorem 1, Corollary 1 in order of appearance, whereas the article prints the same items with the numbering of its own full document.  The complete native source of the article as distributed in the arXiv e-print for this exact version is bundled unchanged as \texttt{sources/209976/source0.tex}.
\begin{theorem}\label{T:2} Let $N$ be a finitely generated $S$-module. The following statements hold true. 
\begin{itemize}			
\item [{\rm (a)}] For each integer $i$, there is an isomorphism $H^i_{\frak n}(N)\cong  H^i_{\frak m}(N)$ of Artinian $R$-modules. Moreover, if $H^i_{\frak n}(N)\neq 0$ then  
$\dim_S(H^i_{\frak n}(N))=\dim_R(H^i_{\frak m}(N)).$
\item [{\rm (b)}] $N$ is Cohen-Macaulay (resp. generalized Cohen-Macaulay, $1$-Cohen-Macaulay,  sequentially  Cohen-Macaulay, sequentially  generalized Cohen-Macaulay, sequentially  $1$-Cohen-Macaulay)  as an $S$-module if and only if so is $N$ as an $R$-module.
\item [{\rm (c)}]  $R$ is a quotient of a Cohen-Macaulay local ring if and only if so is  $S$.
\end{itemize} 		
\end{theorem}

\begin{corollary} \label{C:2} Let $M$ be a finitely generated $R$-module. Then $R\ltimes M$ is sequentially  Cohen-Macaulay (resp. sequentially  generalized Cohen-Macaulay, sequentially  $1$-Cohen-Macaulay) as a ring  if and only if  so is $R\ltimes M$ as an $R$-module, if and only if so are $R$ and $M$.  
\end{corollary}

\begin{proof}  	Set $n=\dim (R)$, $d=\dim_R(M)$ and $S=R\ltimes M$. Then $d\leq n.$ Note that $S=R\oplus M$ as $R$-modules and $\dim_R(R\oplus M)=n$. Therefore, by Theorem \ref{T:2}(b), it is enough to prove that $R\oplus M$ is a sequentially  Cohen-Macaulay (resp. sequentially  generalized Cohen-Macaulay, sequentially  $1$-Cohen-Macaulay)  $R$-module if and only if so are $R$ and $M$. We  proceed by induction on $n$. If  $n=0$ then $R, M, R\oplus M$ are Cohen-Macaulay, so the result  is clear. 
		
Let $n\geq 1.$ Let $S_1$ be the largest $R$-submodule of $R\oplus M$ of dimension less than $n$.  Let $M_1$ be the largest submodule of $M$  of dimension less than $d$ and  $R_1$ the largest $R$-submodule of $R$  of dimension less than $n$.  If $d=n$, we can check that $S_1=R_1\oplus M_1$, therefore $S/S_1\cong R/R_1\oplus M/M_1$. If $d<n$ then $S_1=R_1\oplus M$, so we get   $S/S_1\cong R/R_1$. As $\dim_R(S_1)<n,$ we get by induction that $S_1$ is a sequentially  Cohen-Macaulay (resp. sequentially  generalized Cohen-Macaulay, sequentially  $1$-Cohen-Macaulay) $R$-module if and only if so are $R_1$ and $M_1$. 
		
If $d<n$ then $S/S_1$ is a Cohen-Macaulay (resp. generalized Cohen-Macaulay, $1$-Cohen-Macaulay) $R$-module if and only if so is  $R/R_1$.  So, the result follows.	
		
Assume that $d=n$. Note that $\dim_R(R/R_1)=\dim_R(M/M_1)=n$. Moreover, $$H_{\frak m}^i(S/S_1)\cong H^i_{\frak m}(R/R_1)\oplus H^i_{\frak m}(M/M_1)$$ for all $i<n$. It follows that $S/S_1$ is a Cohen-Macaulay (resp. generalized Cohen-Macaulay, $1$-Cohen-Macaulay) $R$-module if and only if so are $R/R_1$ and $M/M_1$. Thus, we also get the result in this case. 		
\end{proof}

\end{document}
